\section{Cost targets at fixed motion edits}
\label{supp:motion-choice-cost}
Winner labels penalize small and large mistakes alike. Following the general
distinction between prediction and decision cost~\cite{spo}, we compare
winner classification with squared regression of positive gain and positive
harm relative to OLS4. We also test three four-lead cost blocks rather than
one pooled target. The networks, inputs, seeds, update counts, edited tracks,
and peak displacement from OLS4 are identical. Each edit uses one legal
analytic direction shrunk to half the RMS past step, with at least two legal
choices. This is a forced-choice diagnostic, not an admission policy.

\begin{table*}[t]
\centering
\small
\caption{Cost versus winner prediction at identical edits. Entries are gains
in points of the same CV error mass. The bootstrap resamples 14 source groups,
not individual windows. The full-grid column is a missing-label bound, not
a sampling interval. The development recordings were previously examined.}
\label{tab:motion-choice-cost}
\begin{tabular}{lrrr}
\toprule
Contrast & Observed gain & Observed 95\% CI & Full-grid bounds \\
\midrule
Pooled cost vs. classification & \ChoiceCostPooledRank \\
Lead-block cost vs. pooled cost & \ChoiceCostLeadPooled \\
Pooled cost vs. unedited OLS4 & \ChoiceCostPooledOLS \\
\bottomrule
\end{tabular}
\end{table*}

In nine matched fits (three seeds), pooled cost predicted the exact known
winner less often than classification (57.44\% vs. 58.09\%) but reduced
observed regret (1.372 vs. 1.464 pixels per labeled step). Splitting cost by
lead did not improve the pooled result. The missing-label bounds in
Table~\ref{tab:motion-choice-cost} straddle zero; none of the edited arms
passed the unchanged recording-level risk rule. No independent confirmation
or deployment decision follows from this experiment.
